

\subsection{Code source}
\subsubsection{directions}

Commençant d'abord par définir le type enuméré $(enum)$ afin de gérer les différentes directions cardinales d'une manière simple et claire tout au long du project.

Ensuite implémenter la fonction dir\_to\_string qui représente une direction par la chaîne de caractère qui lui est associée aussi que la fonction dir\_to\_short\_string qui représente une direction avec un caractère.

Enfin la fonction dir\_opposite qui renvoie la direction opposée à celle fournie en entrée en manipulant la structure de contrôle 

\subsubsection{locations}

On définie la strucure position $(location\_t)$ par un couple $(x,y)$ désignant respectivement l'abscisse et l'ordonnée d'une position dans le manoir. Pour implémenter ensuite la fonction location qui permet de générer cette structure de position à partir du couple fourni.

Ensuite on implémente les fonctions nécessaires pour manipuler les positions dans le manoir afin de vérifier la validité d'une position dans le manoir, l'égalité entre deux positions et enfin si deux locations sont adjacentes.

\subsubsection{resources}





\subsection{Tests}